\begin{lemma}
Assume \(\varepsilon_1\ge0\), the eventual-comparisons package
\(\noderef{ParameterHierarchyEventualComparisons}
(\varepsilon,\varepsilon_1,\varepsilon_2,n_0)\), \(n>n_0\), and
\[
k=\noderef{TwoBiteNaturalIndependenceScale}(1+\varepsilon,n).
\]
Suppose also that the sample parameters are the rounded hierarchy parameters
\[
m=\noderef{TwoBiteNaturalReducedVertexCount}(n),
\qquad
p=\noderef{TwoBiteEdgeProbability}\!\left(\frac12,n\right).
\]
Let \(I\) be a \(k\)-set.  Conditioned on
\(\mathcal D=\noderef{FiberAndDegreeEvent}\), the huge vertices satisfy the
following bounds.  For \(H_R=H_I\cap V_R\) and \(H_B=H_I\cap V_B\), the
same-color projected contributions
\[
\sum_{x\in H_R}\binom{|\pi_R(X_x(I))|}{2}_{\mathbb R},
\qquad
\sum_{x\in H_B}\binom{|\pi_B(X_x(I))|}{2}_{\mathbb R}
\]
are each \(\varepsilon_1k^2\)-small, and their sum is also
\(\varepsilon_1k^2\)-small.  The two opposite-color projected contributions
are bounded by
\[
(1+\varepsilon_1)\min\left\{\binom{k-\ell}{2},
\binom{pn}{2}+\binom{k-\ell-pn}{2}\right\}
\]
with \(\ell=|\pi_R(I)|\) or \(|\pi_B(I)|\), respectively.
\end{lemma}
\begin{proof}
Let
\[
H_I=\{x:\noderef{TwoBiteHugeClass}(\omega,I,x)\}
\]
be the full huge filter.  Define the two local huge-color filters by
\[
H_R=H_I\cap V_R,\qquad H_B=H_I\cap V_B.
\]
Equivalently, in the Lean statement
\[
\mathsf{hugeRed}(x)\iff
\noderef{TwoBiteHugeClass}(\omega,I,x)\wedge
  \exists r,\ x=\operatorname{inl}(r),
\]
and
\[
\mathsf{hugeBlue}(x)\iff
\noderef{TwoBiteHugeClass}(\omega,I,x)\wedge
  \exists b,\ x=\operatorname{inr}(b).
\]
The finite filters for these predicates are precisely
\[
H_R=\operatorname{univ}.\operatorname{filter}(\mathsf{hugeRed}),
\qquad
H_B=\operatorname{univ}.\operatorname{filter}(\mathsf{hugeBlue}),
\]
which are the Lean filters used in the two child statements and in the
present statement.

We instantiate both child lemmas with the same data as in the present node:
the same sample \(\omega\), the same set \(I\), the same integers \(n\) and
\(m\), the same real \(p\), the same integer \(k\), the same parameters
\(\varepsilon,\varepsilon_1,\varepsilon_2,n_0\), the same nonnegativity
hypothesis \(0\le\varepsilon_1\), the same
\(\noderef{ParameterHierarchyEventualComparisons}\) hypothesis, the same
strict inequality \(n_0<n\), the same rounded sample-parameter identities
\[
m=\noderef{TwoBiteNaturalReducedVertexCount}(n),\qquad
p=\noderef{TwoBiteEdgeProbability}\!\left(\frac12,n\right),
\]
the same rounded independence-scale identity
\[
k=\noderef{TwoBiteNaturalIndependenceScale}(1+\varepsilon,n),
\]
the same \(k\)-set hypothesis \(|I|=k\), and the same event
\(\noderef{FiberAndDegreeEvent}(\omega,\varepsilon_2)\).  Thus every
hierarchy comparison involving \(p,m,k\) is applied only after the same
rounded-parameter identifications as in this node.

First apply \(\noderef{HugeSameColorProjectionBound}\) with this
instantiation.  Set
\[
A_R=\noderef{TwoBiteProjectedPairSum}
  (\omega,I,\mathsf{hugeRed},\noderef{TwoBiteRedProjection}(\omega)),
\]
\[
A_B=\noderef{TwoBiteProjectedPairSum}
  (\omega,I,\mathsf{hugeBlue},\noderef{TwoBiteBlueProjection}(\omega)).
\]
Unfolding \(\noderef{TwoBiteProjectedPairSum}\), these are the finite sums
\[
\sum_{x\in H_R}\noderef{RealChooseTwo}
  \Bigl(\bigl|X_x(I).\operatorname{image}
    (\noderef{TwoBiteRedProjection}(\omega))\bigr|:\mathbb R\Bigr)=A_R,
\]
and
\[
\sum_{x\in H_B}\noderef{RealChooseTwo}
  \Bigl(\bigl|X_x(I).\operatorname{image}
    (\noderef{TwoBiteBlueProjection}(\omega))\bigr|:\mathbb R\Bigr)=A_B.
\]
The displayed cardinalities are the real casts of the finite image
cardinalities in the Lean terms.  In the shorter paper notation these are
\(\sum_{x\in H_R}\binom{|\pi_R(X_x(I))|}{2}_{\mathbb R}\) and
\(\sum_{x\in H_B}\binom{|\pi_B(X_x(I))|}{2}_{\mathbb R}\).
The child conclusion is the triple
\[
\noderef{ClosedPairsBound}(A_R,\varepsilon_1,(k:\mathbb R))
\quad\wedge\quad
\noderef{ClosedPairsBound}(A_B,\varepsilon_1,(k:\mathbb R))
\quad\wedge\quad
\noderef{ClosedPairsBound}(A_R+A_B,\varepsilon_1,(k:\mathbb R)).
\]
Since \(\noderef{ClosedPairsBound}(Q,\varepsilon_1,(k:\mathbb R))\) unfolds
to \(Q\le\varepsilon_1(k:\mathbb R)^2\), these are exactly the red
same-color, blue same-color, and combined same-color conjuncts, in that
order, of the present Lean statement.
Call these three conclusions \(h_{RR}\), \(h_{BB}\), and \(h_{\mathrm{same}}\).
No estimate is extracted from the proof of the child node here: after the
shared hypotheses are supplied, these are simply the three fields of
\(\noderef{HugeSameColorProjectionBound}\), with the same let-bound
\(\mathsf{hugeRed}\) and \(\mathsf{hugeBlue}\) predicates as in the present
statement.

Now apply \(\noderef{HugeOppositeColorProjectionBound}\), again with the same
\(\omega,I,n,m,k,p,\varepsilon,\varepsilon_1,\varepsilon_2,n_0\), the same
\(0\le\varepsilon_1\), the same
\(\noderef{ParameterHierarchyEventualComparisons}\), the same \(n_0<n\), the
same rounded identities for \(m,p,k\), the same \(k\)-set hypothesis, and the
same \(\noderef{FiberAndDegreeEvent}(\omega,\varepsilon_2)\).  Define
\[
C_R=\noderef{TwoBiteProjectedPairSum}
  (\omega,I,\mathsf{hugeRed},\noderef{TwoBiteBlueProjection}(\omega)),
\]
\[
C_B=\noderef{TwoBiteProjectedPairSum}
  (\omega,I,\mathsf{hugeBlue},\noderef{TwoBiteRedProjection}(\omega)).
\]
For the first opposite-color conclusion, the child uses the red huge filter
\(H_R\) but the opposite projection \(\pi_B\).  Unfolding
\(\noderef{TwoBiteProjectedPairSum}\), its left side is
\[
\sum_{x\in H_R}\noderef{RealChooseTwo}
  \Bigl(\bigl|X_x(I).\operatorname{image}
    (\noderef{TwoBiteBlueProjection}(\omega))\bigr|:\mathbb R\Bigr)=C_R,
\]
that is, \(\sum_{x\in H_R}\binom{|\pi_B(X_x(I))|}{2}_{\mathbb R}\).  The
projection-size parameter on the right hand side is not a blue projection
size.  It is
\[
\ell_R=\bigl((I.\operatorname{image}
  (\noderef{TwoBiteRedProjection}(\omega))).\operatorname{card}:\mathbb R\bigr),
\]
the real cast of the red projection image cardinality.  Therefore the child
gives exactly
\[
\begin{aligned}
C_R
&\le
(1+\varepsilon_1)\,
\min\Bigl(
  \noderef{RealChooseTwo}\bigl((k:\mathbb R)-\ell_R\bigr),\\
&\hspace{7em}
  \noderef{RealChooseTwo}\bigl(p\cdot(n:\mathbb R)\bigr)
  +\noderef{RealChooseTwo}
    \bigl((k:\mathbb R)-\ell_R-p\cdot(n:\mathbb R)\bigr)
\Bigr).
\end{aligned}
\]
This is the fourth conjunct of the present Lean statement: its left side is
\[
\noderef{TwoBiteProjectedPairSum}
  (\omega,I,\mathsf{hugeRed},\noderef{TwoBiteBlueProjection}(\omega)),
\]
while its deficit uses the red projection size
\(((I.\operatorname{image}(\noderef{TwoBiteRedProjection}(\omega))).\operatorname{card}:\mathbb R)\),
and its rounded edge-density term is exactly \(p\cdot(n:\mathbb R)\).
Denote this fourth target conjunct by \(h_{RB}\).  The notation is meant to
record the bookkeeping: the huge vertices are red and the summand uses the
opposite blue projection, but the deficit term is measured by the red
projection size of \(I\).

For the second opposite-color conclusion, the child swaps colors: it uses the
blue huge filter \(H_B\) and the opposite projection \(\pi_R\).  Its left side
unfolds to
\[
\sum_{x\in H_B}\noderef{RealChooseTwo}
  \Bigl(\bigl|X_x(I).\operatorname{image}
    (\noderef{TwoBiteRedProjection}(\omega))\bigr|:\mathbb R\Bigr)=C_B,
\]
that is, \(\sum_{x\in H_B}\binom{|\pi_R(X_x(I))|}{2}_{\mathbb R}\).  Its
right-hand projection-size parameter is
\[
\ell_B=\bigl((I.\operatorname{image}
  (\noderef{TwoBiteBlueProjection}(\omega))).\operatorname{card}:\mathbb R\bigr),
\]
the real cast of the blue projection image cardinality.  Therefore the child
gives exactly
\[
\begin{aligned}
C_B
&\le
(1+\varepsilon_1)\,
\min\Bigl(
  \noderef{RealChooseTwo}\bigl((k:\mathbb R)-\ell_B\bigr),\\
&\hspace{7em}
  \noderef{RealChooseTwo}\bigl(p\cdot(n:\mathbb R)\bigr)
  +\noderef{RealChooseTwo}
    \bigl((k:\mathbb R)-\ell_B-p\cdot(n:\mathbb R)\bigr)
\Bigr).
\end{aligned}
\]
This is the fifth conjunct of the present Lean statement: its left side is
\[
\noderef{TwoBiteProjectedPairSum}
  (\omega,I,\mathsf{hugeBlue},\noderef{TwoBiteRedProjection}(\omega)),
\]
but its deficit uses the blue projection size
\(((I.\operatorname{image}(\noderef{TwoBiteBlueProjection}(\omega))).\operatorname{card}:\mathbb R)\),
and again the rounded edge-density term is exactly \(p\cdot(n:\mathbb R)\).
Denote this fifth target conjunct by \(h_{BR}\): the huge vertices are blue,
the summand uses the opposite red projection, and the deficit is measured by
the blue projection size of \(I\).

The three same-color conjuncts from
\(\noderef{HugeSameColorProjectionBound}\) and the two opposite-color
conjuncts from \(\noderef{HugeOppositeColorProjectionBound}\) are, in the
Lean statement order,
\[
h_{RR},\qquad h_{BB},\qquad h_{\mathrm{same}},\qquad h_{RB},\qquad h_{BR}.
\]
The target conjunction is right-associated as
\[
h_{RR}\wedge
  \bigl(h_{BB}\wedge
    (h_{\mathrm{same}}\wedge(h_{RB}\wedge h_{BR}))\bigr),
\]
which is obtained by repeated use of the constructor for conjunction.  No
disjointness argument or additional numerical estimate is being used in this
aggregation step; all numerical work is contained in the two child nodes.
This conjunction assembly proves the lemma.
\end{proof}
